\section*{The actual entire radial trace of the orthogonal swirl component}

\subsection*{A full state map retaining the spatial complement and every radial input}

Retain the physical Hilbert spaces and operators of the solenoidal
correspondence: \(H=L^2(\mathbb R,ds;\mathbb C^3)\),
\(V=L^2(\mathbb R^3,dx;\mathbb C^3)\), its solenoidal subspace
\(V_\sigma\), \(K_\nu\), \(P_\nu=K_\nu K_\nu^*\),
\(\Pi u=\langle r u_\theta\rangle\), and the section \(\mathcal R\).
Here \(\sigma=r^2/2\) is unchanged. Write
\(C_{\rm sw}=\mathcal R\Pi\) for the orthogonal operator previously
called \(C\); the source amplitude \(C>1\) remains a positive scalar.
This typed renaming leaves both objects unchanged. Put
\[
 H_\infty=\bigcap_{m\geq0}\operatorname{Dom} A^m,
 \qquad V_\infty=\bigcap_{m\geq0}\operatorname{Dom}B_\nu^m,
 \qquad V_{{\rm rem},\infty}
 =V_\infty\cap V_\sigma\cap\ker K_\nu^*\cap\ker\Pi,
\]
where \(Ah=s^2h/4\) and \(B_\nu=-\nu\Delta\) have their full
Fourier multiplication domains. Let
\[
 \mathscr B_c=\{B\in C_c^\infty([0,\infty)\times\mathbb R):
                                      B(0,z)=0\text{ for every }z\},
 \qquad
 \mathscr U=\{u\in V_\infty\cap V_\sigma:\Pi u\in\mathscr B_c\}.
\]
The actual complete corrected source belongs to \(\mathscr U\) at each
preterminal time, by its compact support and the proved angular endpoint
orders. For every \(B\in\mathscr B_c\),
\(B/\sigma=\int_0^1B_\sigma(a\sigma,z)\,da\) has its smooth
extension at zero. Thus \(\mathcal RB\) is a compactly supported
smooth solenoidal vector field and lies in every displayed domain.
In particular its weighted energy is finite, with the unchanged measure
\(2\pi(2\sigma)^{-1}d\sigma\,dz\).

The complete orthogonal correspondence already proved restricts to the
following bijection with the displayed full target:
\[
 \mathscr U\longleftrightarrow
 H_\infty\oplus\mathscr B_c\oplus V_{{\rm rem},\infty}.
\]
\[
 \begin{gathered}
 u\longmapsto(h_{\rm sol},B,v)
   =(K_\nu^*u,\Pi u,(I-P_\nu-C_{\rm sw})u),\\
 u=K_\nu h_{\rm sol}+\mathcal RB+v.
 \end{gathered}
\]
Indeed preservation of every generator domain proves the forward
inclusions. Conversely each displayed triple reconstructs a field in
\(V_\infty\cap V_\sigma\); applying \(\Pi\) returns \(B\), because
\(\Pi K_\nu=0\), \(\Pi\mathcal R=I\), and \(\Pi v=0\).
The identities \(K_\nu^*K_\nu=I\), \(K_\nu^*\mathcal R=0\) and
\(K_\nu^*v=0\) return \(h_{\rm sol}\); subtraction returns \(v\).
These prove both inverse identities and surjectivity. Real fields retain
the original positive-real/negative-pure-imaginary branch condition in
\(h_{\rm sol}\), with real \(B,v\).

For arbitrary \(B\in\mathscr B_c\), use the complete sequence ingress
\(\mathfrak I\) of Section~\ref{sec:cq-angular-ingress}, with its one
fixed cutoff and \(\epsilon_n=2^{-n}\). Define
\[
 \mathscr P(B)=\bigl(\mathfrak I(B(\cdot,z))\bigr)_{z\in\mathbb R}.
\]
Every entry is an actual order-one entire radial trace. Its precise
target is the set of such families whose individually recovered inputs
fit a function in \(\mathscr B_c\), with the fixed cutoff sequence
specified there. This is a checked image condition on the explicit
inverse, rather than a claim about arbitrary families of entire functions.
For \(\sigma>0\), inversion of the \(n\)-th entry gives exactly
\(\chi(\sigma/\epsilon_n)B(\sigma,z)\). Once
\(2\epsilon_n<\sigma\), it equals \(B\) on a neighborhood of that
point. The smooth recovered function supplies every endpoint derivative;
its value at zero is the retained value zero. Thus \(\mathscr P\) is
injective and its stated inverse recovers the full joint input.
Consequently the full state has the equally exact encoding
\[
 u\longleftrightarrow(h_{\rm sol},\mathscr P(B),v),
 \qquad
 u=K_\nu h_{\rm sol}+\mathcal R\mathscr P^{-1}(\mathscr P(B))+v.
\]
No finite sequence entry or axial restriction replaces the full input.
Its energy remains
\[
 \|u\|_V^2=\|h_{\rm sol}\|_H^2+
 2\pi\int_{\mathbb R}\int_0^\infty
                \frac{|B(\sigma,z)|^2}{2\sigma}\,d\sigma\,dz
 +\|v\|_V^2.
\]

\subsection*{The central entire observable and its exact axis functional}

On the actual central slab \(|z|<\sqrt\nu z_0/2\), the complete
coefficient recurrence, all six complex Picard rows, every cutoff term,
and the full physical force in
Section~\ref{sec:cq-actual-angular-analytic} prove endpoint analyticity
of the untruncated \(B,P,Q,F\) at every \(0\leq t<1\).
The single entire observable there is exactly
\[
 \mathcal T_{h,z}(u)(s)=T_h(\Pi u)(s,z)
 =K_h(w)\widehat{\Pi u}(w/2,z),\qquad w=\tfrac12-is,
\]
on the subspace of \(\mathscr U\) whose \(\Pi u(\cdot,z)\) has an
analytic endpoint germ. Its target lies in \((\cqE,\gamma)\), and
the inverse of \(T_h\) is precisely the full radial Mellin and two
dyadic inverse series already proved. In particular
\[
 \mathcal T_{h,z}(u)=\mathcal T_{h,z}(\mathcal R\Pi u),
 \qquad
 \ker\mathcal T_{h,z}=
 \{u\text{ in this domain}:\Pi u(\cdot,z)=0\}.
\]
The first equality follows by \(\Pi\mathcal R\Pi=\Pi\); the second
follows by the injectivity of \(T_h\). Both \(K_\nu h_{\rm sol}\)
and \(v\) belong to that kernel whenever they arise from the full
state above. This specifies exactly what this one axial observable sees;
the full encoding \((h_{\rm sol},\mathscr P(B),v)\) still recovers
all of \(u\).

There is an exact numerical correspondence between the earlier direct
Mellin pole and the value of this entire function. For any endpoint
analytic \(a\in C_c^\infty([0,\infty))\) with \(a(0)=0\), define
\(s_0=-5i/2\). The full radial residue formula and the proved value
\(\zeta'(-2)=-\zeta(3)/(4\pi^2)\) give
\begin{align*}
 (T_ha)(s_0)
 &=\frac{\zeta(3)}{8\pi^2}
       (1-2^{-3/2})(2^{h+1}-1)a'(0),\\
 \frac{(T_ha)(s_0)}{\Xi(s_0)}
 &=\frac{(1-2^{-3/2})(2^{h+1}-1)}{12\pi}a'(0),
 \qquad \Xi(s_0)=\frac{3\zeta(3)}{2\pi}.
\end{align*}
Indeed the residue of \(M_a^{\rm rad}\) at \(w=-2\) is
\(2^{-1}(1-2^{-3/2})(1-2^{h+1})a'(0)\). Multiplication by the
displayed negative derivative of zeta gives the first formula with all
signs and factors; division by the nonzero displayed \(\Xi(s_0)\)
gives the second. The direct arithmetic endpoint calculation gives
\((\mathcal A_ha)'(0)=(2^{h+1}-1)a'(0)/16\), hence
\[
 (T_ha)(s_0)=\frac{2\zeta(3)}{\pi^2}(1-2^{-3/2})
                    (\mathcal A_ha)'(0),
 \qquad
 \frac{(T_ha)(s_0)}{\Xi(s_0)}
   =\frac4{3\pi}(1-2^{-3/2})(\mathcal A_ha)'(0).
\]
All scalar coefficients are nonzero. Thus these are invertible scalar
maps between the endpoint derivative, the direct Mellin residue, and the
fixed entire value on this input class. They do not identify the two
full transforms; their distinct arguments and full inverses remain
those proved in the imported radial correspondence.

For the actual field, the independently proved circulation identity
\(\partial_\sigma\Pi u(0,z)=\omega_3(0,0,z)\) now yields
\[
 \mathcal T_{h,z}(u)(s_0)
 =\frac{\zeta(3)}{8\pi^2}(1-2^{-3/2})(2^{h+1}-1)
                          \omega_3(0,0,z).
\]
Both other orthogonal summands have zero value of this functional.
At \(z=0,t=1-\tau\), the actual coefficient
\(\omega_{3,\nu}(0,0,0,1-\tau)=2\tau^{-1-h}/C\) therefore proves
\[
 \mathcal T_{h,0}(u_\nu)(s_0,1-\tau)
 =\mathcal T_{h,0}(\mathcal R B_\nu)(s_0,1-\tau)
 =\frac{\zeta(3)(1-2^{-3/2})(2^{h+1}-1)}{4\pi^2C}
                                                   \tau^{-1-h}.
\]
The interval is exactly \(0<\tau<\tau_1\) for this axis identity;
the source amplitude, exponent and viscosity transformation are unchanged.

\subsection*{The retained nonlinear evolution and source quotient}

The actual four central traces obey the entire angular equation
\eqref{eq:cq-radial-entire-angular-pde}, with its exact residual
\(\kappa_h(s-2i)\mathcal F_E^{\rm rad}(s)\).
In the present full state, write \(u^{(1)}=K_\nu h_{\rm sol}\),
\(u^{(2)}=\mathcal R B\), and \(u^{(3)}=v\). Its actual fluxes are
still precisely
\[
 P=\sum_{a,b=1}^3\langle r^2u_r^{(a)}u_\theta^{(b)}\rangle,
 \qquad
 Q=\sum_{a,b=1}^3\langle r u_z^{(a)}u_\theta^{(b)}\rangle.
\]
These follow by expanding both factors in the original definitions.
All nine ordered terms are retained, even though \(\Pi u^{(1)}\)
and \(\Pi u^{(3)}\) vanish. Only their complete sums are used as
the compact endpoint-analytic inputs to \(T_h\); no compact-support
or analyticity assertion is imposed on an individual nonlocal summand.
The original pressure equation remains the longitudinal projection of
\(f-\sum_{a,b}N(u^{(a)},u^{(b)})\), with its spatial constant part
retained. The source force has every cutoff and physical viscosity term
already displayed. The full inverse state map and the injective residual
calculation consequently recover the same physical angular equation,
together with the other two component equations from the preceding
orthogonal correspondence.

Finally write \(q:\cqE\to\cqQ=\cqE/\cqN\) for the original source
projection. Its composite with this actual trace has the exact kernel
\[
 \ker(q\mathcal T_{h,z})
 =\{u\text{ in the stated domain}:T_h(\Pi u(\cdot,z))\in\cqN\}.
\]
The inverse before this projection remains the radial inverse above.
At the fixed point \(s_0\), put \(e=\Xi/\Xi(s_0)\). The original
polynomial trace subspace gives \(e\in\cqN\) and \(e(s_0)=1\).
The representative relation has the exact image
\[
 \{(q(f),f(s_0)):f\in\cqE\}=\cqQ\times\mathbb C.
\]
To prove surjectivity, choose any representative \(f\) of the prescribed
class and replace it by \(f+(b-f(s_0))e\); its class is unchanged and
its value is exactly the prescribed \(b\). Its kernel as a linear map
from \(\cqE\) to that product is exactly
\(\cqN\cap\{f:f(s_0)=0\}\), by the two definitions.
This gives the complete algebraic correspondence for the loss of this
representative value. No continuity into the unnamed source topology,
identification of that topology as Fr\'echet, ideal property of the
actual \(\cqN\), or membership of \(\Xi'\) is inferred.
